%%%%%%%%%%%%%%%%%%%%%%%%%%%%%%%%%%%%%%%%
\chapter{Introducción}
%%%%%%%%%%%%%%%%%%%%%%%%%%%%%%%%%%%%%%%%


El aprendizaje profundo se ha consolidado como una de las líneas más influyentes de la inteligencia artificial aplicada al análisis de imágenes. Su desarrollo reciente ha estado acompañado por avances en capacidad computacional, disponibilidad de datos y diseño de arquitecturas cada vez más eficaces para procesar información visual de alta complejidad \parencite{herath2025evolution,mienie2024review}. En este marco, la visión por computador ha experimentado una expansión sostenida, especialmente en tareas como clasificación, detección y segmentación, en las que los modelos pueden aprender representaciones jerárquicas a partir de los datos sin necesidad de diseñar manualmente todas las características relevantes.

Dentro de este campo, las redes neuronales convolucionales han ocupado un lugar central. Su importancia radica en que permiten capturar patrones espaciales, reconocer estructuras visuales locales y combinarlas progresivamente en representaciones más abstractas \parencite{khan2020survey,raj2025cnn}. Gracias a ello, las redes neuronales convolucionales (\textit{Convolutional Neural Networks}, CNN) y sus variantes se han consolidado como una base técnica dominante en numerosos problemas de reconocimiento visual, especialmente cuando los objetos de interés presentan variaciones de forma, tamaño, posición, iluminación o fondo.

La relevancia de estas arquitecturas no se limita a contextos industriales o comerciales. También ha alcanzado dominios científicos en los que el reconocimiento visual constituye una tarea clave para la generación de conocimiento. Entre ellos se encuentra el estudio de la biodiversidad, ámbito en el que la identificación de organismos a partir de imágenes ha ganado importancia como apoyo al monitoreo ecológico, la documentación de especies y la observación sistemática de cambios en comunidades biológicas \parencite{borowiec2022deep,hoye2021entomology}. En este escenario, el uso de aprendizaje profundo representa una oportunidad para fortalecer procesos de análisis que tradicionalmente han dependido de trabajo manual intensivo.

Uno de los casos más relevantes dentro de esta intersección entre visión por computador y ecología es el reconocimiento de polinizadores a partir de imágenes. Los polinizadores, particularmente abejas, abejorros y otros insectos asociados a la polinización, cumplen una función esencial en la reproducción de numerosas plantas y en el mantenimiento de la diversidad biológica \parencite{ghosh2025pollinator}. Por esta razón, su observación e identificación resultan valiosas para el seguimiento de comunidades ecológicas, la evaluación de cambios en poblaciones y la comprensión de procesos vinculados con la conservación \parencite{hellwig2024pollinator}.

A pesar de esta importancia, la identificación de polinizadores sigue presentando dificultades prácticas y metodológicas. Gran parte del reconocimiento de especies continúa dependiendo de observación experta, análisis morfológico y revisión manual de imágenes. Estos enfoques requieren tiempo, entrenamiento especializado y condiciones de captura relativamente favorables para distinguir organismos que, en muchos casos, presentan una alta similitud visual entre sí \parencite{spiesman2021bumblebee}. Además, la variabilidad de iluminación, postura, escala, movimiento y fondo puede dificultar aún más la diferenciación entre especies o grupos taxonómicos \parencite{ratnayake2021agriculture}.

Estas limitaciones hacen evidente la necesidad de contar con herramientas computacionales capaces de asistir o automatizar parte del proceso de reconocimiento. En este punto, el aprendizaje profundo aparece como una alternativa especialmente prometedora. En particular, las redes neuronales convolucionales permiten aprender directamente a partir de imágenes rasgos visuales que serían difíciles de formalizar mediante reglas manuales. Esta capacidad resulta especialmente relevante cuando se trabaja con organismos pequeños, con alta variación intraespecífica o con entornos visuales complejos, características frecuentes en el reconocimiento de polinizadores \parencite{hoye2021entomology,borowiec2022deep}.

El potencial de estas técnicas ya se refleja en estudios recientes que reportan aplicaciones de aprendizaje profundo para clasificación de especies, detección de insectos en imágenes y monitoreo automatizado en tiempo real \parencite{bjerge2022tracking,stark2023yolo}. Asimismo, se han propuesto sistemas orientados a contextos específicos de observación, lo que evidencia que el reconocimiento automatizado de polinizadores mediante imágenes ha dejado de ser una posibilidad teórica para convertirse en una línea activa de investigación aplicada \parencite{urbano2025embedded,stark2025pollinators}. Sin embargo, este crecimiento no ha ido acompañado de una síntesis suficientemente clara que permita comparar enfoques, métricas y decisiones técnicas dentro del campo.

Al revisar la literatura disponible, se observa que los avances no se presentan de forma homogénea ni fácilmente comparable. Los estudios difieren en los grupos de polinizadores analizados, en el tipo de imágenes empleadas, en las tareas abordadas y en las arquitecturas implementadas. Mientras algunos trabajos se concentran en clasificación, otros priorizan detección, seguimiento automatizado o segmentación; además, las condiciones de captura, el tamaño de los conjuntos de datos y el nivel de detalle taxonómico también varían de manera importante \parencite{spiesman2021bumblebee,proenca2025leafhoppers}. Los estudios son heterogéneos, por lo que requieren una síntesis organizada que permita compararlos dentro de grupos metodológicamente semejantes.

La comparación entre estudios también se ve limitada por la diversidad de métricas reportadas y por la descripción desigual de las decisiones técnicas que acompañan el entrenamiento de los modelos. No todos los trabajos informan las mismas medidas de desempeño, ni lo hacen con el mismo nivel de detalle; además, factores como la disponibilidad limitada de datos etiquetados, el desbalance entre clases, la variabilidad morfológica intra e interespecífica, la transferencia de aprendizaje, el aumento de datos y los métodos de optimización pueden influir de manera importante en los resultados \parencite{khanzhina2022combating}. 

En este contexto, la ausencia de una síntesis estructurada constituye un problema académico concreto, ya que dificulta responder con claridad qué arquitecturas se utilizan con mayor frecuencia, qué métricas predominan y qué factores técnicos se asocian con mejores niveles de desempeño reportado.

La pertinencia de esta investigación radica en la necesidad de organizar y comparar una literatura reciente que, aunque valiosa, permanece dispersa. Una revisión sistemática permite reunir estudios bajo criterios explícitos de selección, sintetizar sus enfoques y ofrecer una visión integrada de los componentes más relevantes del campo. En áreas afines, este tipo de revisión ha resultado útil para clarificar el desarrollo de arquitecturas, patrones de evaluación y desafíos persistentes \parencite{pacal2024systematic,wang2025plantreview}. 

En el caso del reconocimiento de polinizadores mediante aprendizaje profundo, una síntesis de esta naturaleza puede servir para describir el estado del arte y orientar futuras decisiones metodológicas en torno a la selección de modelos, métricas y estrategias técnicas.

Desde esta perspectiva, la presente monografía tiene como propósito caracterizar el estado del arte del reconocimiento de imágenes de polinizadores mediante aprendizaje profundo, a partir de estudios científicos publicados desde 2020. El análisis se centra en tareas de clasificación, detección y segmentación, con énfasis en el papel de las redes neuronales convolucionales y de otras arquitecturas relacionadas, por constituir uno de los enfoques más representativos en problemas de reconocimiento visual.

De manera más específica, el estudio busca identificar las arquitecturas y tipos de modelos más utilizados, sintetizar las métricas de evaluación reportadas en la literatura y analizar los principales factores técnicos asociados al rendimiento de los modelos. En conjunto, estas dimensiones permiten organizar la evidencia reciente desde una perspectiva comparativa y ofrecer una visión integrada de un campo que ha crecido de manera notable, pero cuya consolidación metodológica aún presenta limitaciones importantes.
